\scanpage{75}
% archive sheet number: 74

Les développements précédents étaient centrés sur le groupe $H=\unsure{\mathrm{Gp}(1,\mathbb C)|\mathrm{GP}(1,\mathbb C)}$ opérant sur $\mathbb P^1(\mathbb C)$ (ou $\mathrm{Gp}(1)_{\mathbb Z}$ sur $\mathbb P^1_{\mathbb Z}\ldots$), et les \unsure{invariants|invariant} différentiels, sous l'action de ce groupe, des fonctions analytiques complexes (\uline{NB} marcherait aussi dans l'analytique réel, ou les $C^\infty$, voire $C^3$). On va passer rapidement en revue les développements analogues pour d'autres groupes

a) \uline{Groupe des translations} $\mathbb G_a$ \uline{opérant sur} $\mathbb E^1$

On a l'opération (dans l'analytique complexe p.ex)

\[
d:\underline{\mathcal O}_X\longrightarrow\underline\Omega^1_{X/\mathbb C}=\underline\omega_X
\]

% slanted
\margin{\uline{Ici} $\underline{\mathcal O}_X$ mauvais pt de vue, prendre $\unsure{\mathcal E_X|\mathscr E_X}$ plutôt !}

qui est \inserted{épi ($X$ lisse dim 1) et} nul \deleted{exactement} sur les fonctions loc. constantes, donc posant

\[
A(f,g)=dg-df,\qquad A:\underline{\mathcal O}_X\times\underline{\mathcal O}_X\longrightarrow\underline\omega_X
\]

on trouve l'analogue de a) b) dessus, plus

\[
 A(f,g+c)=A(f,g)\qquad c\in\mathbb C
\]

plus

\[
 A(f,g)=A(f,h),\quad g(x)=h(x)\implies \underset{\text{au vois. de }x}{g=h}
\]

et

\[
g\mapsto A(f,g)\qquad \underline{\mathcal O}_X\longrightarrow\underline\omega_X\quad \text{épi}
\]

.(Cas $X$ lisse de dim 1 pr. \unsure{st. mt|smt|définition}) --

Ici est pas la peine d'introduire de torseur, \inserted{(si on veut, il est trivial)} ni $=$ $A(f,g)$, car l'expression $dg$ contient \unsure{toutes les|tous les} qui en sortent
